% Generated draft: requires mathematical and source review.
\documentclass[11pt]{article}
\usepackage[margin=25mm]{geometry}
\usepackage{fontspec}
\setmainfont{Latin Modern Roman}
\usepackage{amsmath,amssymb,mathtools,unicode-math}
\setmathfont{Latin Modern Math}
\usepackage{enumitem,needspace,xurl,hyperref}
\setlength{\parindent}{0pt}
\setlength{\parskip}{5pt}
\setlength{\emergencystretch}{3em}
\newcommand{\N}{\mathbb{N}}
\newcommand{\Z}{\mathbb{Z}}
\newcommand{\Q}{\mathbb{Q}}
\newcommand{\R}{\mathbb{R}}
\newcommand{\C}{\mathbb{C}}
\newcommand{\F}{\mathbb{F}}
\newcommand{\A}{\mathbb{A}}
\newcommand{\PP}{\mathbb{P}}
\newcommand{\E}{\mathbb{E}}
\newcommand{\eps}{\varepsilon}
\newcommand{\dd}{\,\mathrm{d}}
\begin{document}
\section*{UnsolvedMath ID: 30003452}
\textbf{Source problem code:} catalogue entry
\subsection{Definitions and mathematical statement}
For terminology, under the standard complex monic-pencil convention, let \(g,d\geq 1\) and let \(A=(A_1,\ldots,A_g)\in M_d(\C)^g\). For \(n\geq 1\) and \(X=(X_1,\ldots,X_g)\in M_n(\C)^g\), set
\[
L_A(X)=I_d\otimes I_n+\sum_{j=1}^g\bigl(A_j\otimes X_j+A_j^*\otimes X_j^*\bigr),
\qquad
\mathcal D_A(n)=\{X\in M_n(\C)^g:L_A(X)\succeq 0\}.
\]
Here \(^*\) denotes conjugate transpose and \(\succeq 0\) denotes positive semidefiniteness. The sequence \(\mathcal D_A=(\mathcal D_A(n))_{n\geq 1}\) is the free spectrahedron determined by \(A\). In this convention, boundedness means that there is a constant \(R<\infty\) such that \(\|X_j\|\leq R\) for every level \(n\), every \(X\in\mathcal D_A(n)\), and every \(1\leq j\leq g\). Define \(\mathcal D_B\) analogously for a coefficient tuple \(B\), whose matrix size may differ from that of \(A\).

A free analytic map is a level-preserving map that is holomorphic at each matrix level and respects direct sums and simultaneous similarities whenever the tuples involved belong to its domain. A free bianalytic map is a bijective free analytic map whose inverse is free analytic. Analytic maps between free spectrahedra are ordinarily considered on their levelwise interiors. These definitions provide terminology only: the supplied record does not specify the underlying pencil convention, scalar field, tuple lengths, or whether the referenced theorem imposes additional boundary conditions.

\textbf{Source question, with unresolved theorem reference.} In the partial converse classifying bianalytic maps between bounded free spectrahedra \(\mathcal D_A\) and \(\mathcal D_B\), can the genericity hypothesis on the coefficient tuples \(A\) and \(B\) be removed while retaining every other hypothesis and the entire conclusion of that partial converse?

More explicitly, if the referenced partial converse has the schematic form
\[
H(A,B,f)\ \land\ G(A,B)\ \Longrightarrow\ C(A,B,f),
\]
the question asks whether
\[
H(A,B,f)\ \Longrightarrow\ C(A,B,f)
\]
holds for every pair of coefficient tuples and every bianalytic map covered by its remaining hypotheses. Here \(H\), \(G\), and \(C\) are placeholders for, respectively, all remaining hypotheses, the particular genericity hypothesis, and the precise classification conclusion of the referenced theorem. The supplied record does not reproduce these ingredients or identify a theorem number. Consequently, this is a qualified formulation of the supplied question, not a fully specified classification conjecture; the missing hypotheses and conclusion must be recovered from the cited source before a precise standalone theorem statement can be given.
\subsection{Short English statement}
The question asks whether the partial converse classifying bianalytic maps between bounded free spectrahedra remains valid without assuming that the coefficient tuples \(A\) and \(B\) are generic. Every other hypothesis and the classification conclusion must be retained. The supplied record omits the exact partial converse and its definition of genericity, so the intended universal statement remains incompletely specified.
\subsection{Sources}
\begin{itemize}
\item Supplied local UnsolvedMath registry record, permanent ID \(30003452\), title ``Bianalytic Free Spectrahedron Maps Without Genericity,'' source problem code \texttt{catalogue entry}, in the Oberwolfach Reports Open Problems collection. The record gives proposed year \(2017\) and names no proposer.
\item Original Oberwolfach report: \emph{Real Algebraic Geometry With a View Toward Moment Problems and Optimization} (2018), \emph{Oberwolfach Reports}, pp.\ 771--862, DOI \texttt{10.4171/owr/2017/14}. Supplied source and literature link: \url{https://doi.org/10.4171/owr/2017/14}. The record describes these reports as documenting mathematical workshops and their open-problem discussions.
\item Supplied literature reference: \emph{Bianalytic maps between free spectrahedra}, \url{https://doi.org/10.1007/s00208-017-1630-3}.
\item Supplied literature reference: \emph{Free bianalytic maps in a generic setting}, \url{https://arxiv.org/abs/1711.09459}.
\item Literature assessment dated exactly \texttt{2026-08-21}: the supplied curation labels the question open and states that published results still require genericity or related irreducibility assumptions, leaving the extracted generalization open. This is supplied dated curation derived from the Oberwolfach material, not an independently verified literature assessment. The record itself requests independent verification of time-sensitive status and cited claims.
\item The registry assigns the statement status \texttt{corrected\_verified}, explaining that notation, typography, and sentence structure were repaired while mathematical content was preserved. This supplied editorial label does not establish mathematical correctness or resolve the missing theorem reference. No cited source was independently consulted for this draft.
\end{itemize}
\end{document}
